%May, 23 2008

%\documentclass[11pt]{article}
\documentclass[graybox]{./styles/svmult}
\usepackage{graphicx}
\usepackage{helvet}
\usepackage{courier}
\usepackage{makeidx}
\usepackage{multicol}
\usepackage{mathptmx}
\usepackage{type1cm}
\usepackage[bottom]{footmisc}
%\usepackage{footmisc}
\usepackage[authoryear]{natbib}

\usepackage{theorem}
\usepackage{times} %greatly improves the on-screen readability

%old packages

\usepackage{optprog}
\usepackage{amssymb,amsmath}
%\usepackage{amsthm}
\usepackage{rotating}
\usepackage{setspace}
\usepackage{lscape}
\usepackage{verbatim}

\newcommand{\zaupdate}[1]{\noindent \textcolor{blue}{\bf ZE -- #1 } }

\newcommand{\fvarj}{t_j}
\newcommand{\colq}{q = ({\mbox{x}^q}^T, {\mbox{y}^q}^T, {\mbox{t}^q}^T )^T}
\newcommand{\zq}{q' = ({\mbox{x}^{q'}}^T, {\mbox{y}^{q'}}^T, {\mbox{t}^{q'}}^T )^T}
\newcommand{\Z}{\mathbb{Z}}

%\setlength{\evensidemargin}{0in}
%\setlength{\oddsidemargin}{0in}
%\setlength{\textwidth}{6.5in}
%\setlength{\parindent}{0in}{\bf }
%\setlength{\parskip}{0.06in}
%\setlength{\textheight}{9.0in}
%\onehalfspacing
\makeindex

\begin{document}
  \title*{A Branch-and-Price Algorithm for Combined Location and Routing Problems Under Capacity Restrictions}
\author{Z. Akca \and R.T. Berger \and T.K. Ralphs}
\institute {Z. Akca \at Department of Industrial and Systems Engineering, Lehigh University, Bethlehem, PA 18015, \\
\email{zea2@lehigh.edu}
\and
R.T. Berger \at 18 Whittemore Street, Concord, MA 01742, \email{rosemary.berger@verizon.net}
\and
T.K. Ralphs \at Department of Industrial and Systems Engineering, Lehigh University, Bethlehem, PA 18015, \\
\email{ted@lehigh.edu}, http://coral.ie.lehigh.edu/\~{ }ted }

\titlerunning{A Branch-and-Price Algorithm for LRPs Under Capacity Restrictions }
\maketitle
\abstract{
We investigate the problem of simultaneously determining
the location of facilities and the design of vehicle routes to
serve customer demands under vehicle and facility capacity
restrictions. We present a set-partitioning-based formulation of
the problem and study the relationship between this formulation
and the graph-based formulations that have been used in previous
studies of this problem. We describe a branch-and-price algorithm
based on the set-partitioning formulation and discuss
computational experience with both exact and heuristic
variants of this algorithm.
\keywords{branch-and-price, facility location, vehicle routing, column generation} 
}
   
\section {Introduction}
The design of a distribution system begins with the questions of where
to locate the facilities and how to allocate customers to the selected
facilities. These questions can be answered
using location-allocation models, which are based on the assumption that
customers are served individually on out-and-back routes. However, when
customers have demands that are
less-than-truckload and thus can receive service from routes
making multiple stops, the assumption of individual routes will not
accurately capture the transportation cost. Therefore, the integration
of location-allocation and routing decisions may yield more accurate
and cost-effective solutions.

In this paper, we investigate the so-called \emph{location and routing
problem} (LRP). Given a set of candidate facility locations and a set
of customer locations, the objective of the LRP is to determine the
number and location of facilities and construct a set of vehicle
routes from facilities to customers in such a way as to minimize total
system cost. The system cost may include both the fixed and operating
costs of both facilities and vehicles. In this study, we consider an
LRP with capacity constraints on both the facilities and the vehicles.

Vehicle routing models in general have been the subject of a wide range of
academic papers, but the number of these devoted to combined location and
routing models is much smaller. \citet{Laporte} surveyed the work on
deterministic LRPs and described different formulations of the problem,
solution methods used, and the computational results published up to 1988.
\citet{Min} classified both deterministic and stochastic models in the LRP
literature with respect to problem characteristics and solution methods. 

Solution approaches for the LRP can be divided broadly into heuristics
and exact methods, with much of the literature devoted to heuristic
approaches. Most heuristic algorithms divide the problem into
subproblems and handle these subproblems sequentially or iteratively.
Examples of heuristics developed for the LRP can be found in
\citet{Perl}, \citet{Hansen} and \citet{Barreto07}. 

Many fewer papers have been devoted to the study of exact algorithms
for the LRP and most of these have been based on the so-called
two-index vehicle flow formulation. \citet{Laporte81} solved a single
depot model with a constraint relaxation method and a branch-and-bound
algorithm. They reported solving problems with 50 customer locations.
\citet{Laporte83} solved a multi-depot model using a constraint
relaxation method and Gomory cutting planes to satisfy integrality.
They were able to solve problems with at most 40 customer sites.
\citet{Laporte86} applied a branch-and-cut algorithm to a multi-depot
LRP model with vehicle capacity constraints. They used subtour
elimination constraints and chain barring constraints that guarantee
that each route starts and ends at the same facility. They reported
computational results for problems with 20 customer locations and 8
depots. Finally, \citet{Belenguer} provided a two-index formulation of
the LRP with capacitated facilities and capacitated vehicles and
presented a set of valid inequalities for the problem. They developed
two branch-and-cut algorithms based on different formulations of the
problem. They reported that they could solve instances with up to
32 customers to optimality in less than 70 CPU seconds and could provide
good lower bounds for the rest of the instances, which had up to 134
customers.

\citet{Berger} developed a set-partitioning-based model for a special
case of the LRP with route length constraints and uncapacitated
vehicles and facilities. \citet{Berger07} extended that work to
develop an exact branch-and-price algorithm in which they solved the
pricing problem as an elementary shortest path problem with one
resource constraint. They reported computational results for problems
with 100 customers and various distance constraints. \citet{Akca}
developed an exact solution algorithm using branch-and-price
methodology for the integrated location routing and scheduling problem
(LRSP), which is a generalization of the LRP. In the LRSP, the
assumption that each vehicle covers exactly one route is removed and
the decision of assigning routes to vehicles subject to the scheduling
constraints is considered in conjunction with the location and routing
decisions. They considered instances with capacitated facilities and
time- and capacity-limited vehicles. They provided solutions for
instances with up to 40 customers.
 
In this study, we utilize a variant of the model presented by
\citet{Akca} to solve the LRP under capacity restrictions and modify
their exact algorithm for the LRSP to solve the LRP. We develop a
number of variants and heuristic extensions of the basic algorithm and
report on our computational experience solving both randomly generated
instances and instances from the literature. The remainder of the
paper is organized as follows. Section \ref{sec:ProbDef} presents a
formal description of the problem, provides two formulations, and
investigates the relationship between the formulations. Section
\ref{sec:SolAlg} describes details of the heuristic and the exact
algorithms for the set-partitioning formulation. Section \ref{sec:CompR}
provides some computational results evaluating the performance of the
algorithms. Section \ref{sec:Conc} concludes the paper.
\section{Problem Definition and Formulations}
\label{sec:ProbDef}
The objective of the LRP is to select a subset of facilities and
construct an associated set of vehicle routes serving customers at
minimum total cost, where the cost includes the fixed and operating
costs of both facilities and vehicles. The constraints of the problem
are as follows: (i) each customer must be serviced by exactly one
vehicle, (ii) each vehicle must be assigned to exactly one facility at
which it must start and end its route, (iii) the total demand of the
customers assigned to a route must not exceed the vehicle capacity,
and (iv) the total demand of the customers assigned to a facility must
not exceed the capacity of the facility. In the literature, most of
the exact methods developed for the described LRP or its special cases
are based on the two-index vehicle flow formulation of the problem. To
the best of our knowledge, an exact solution algorithm based on a
set-partitioning formulation has not been applied to the case of the
LRP with capacity constraints. The theoretical relationship between
the two-index formulation and the set-partitioning formulation can be
understood by considering a closely related three-index formulation
that we present below. 
We show that applying Dantzig-Wolfe
decomposition to the three-index formulation yields the
set-partitioning formulation. This is turn shows that the bounds
yielded by the LP relaxation of the set-partitioning model must be at
least as tight as those of the three-index formulation.

\subsection{Vehicle Flow Formulation}

We let $I$ denote the set of customers, $J$ denote the set of
candidate facilities, and $V = I \cup J$. To bypass the decision 
of assigning vehicles to facilities, we assume that each facility
has its own set of vehicles and that a vehicle located at facility $j$ can
only visit customer locations and the facility $j$ during its trip.
Let $H_j$ be the set of
vehicles located at facility $j$, $\forall j \in J$ and $H$ be the set
of all vehicles, $H = \bigcup_{j \in J}H_j$. We define the following
parameters and decision variables:

\noindent\underline{Parameters:}
\begin{eqnarray*}
D_i & = & \mbox{demand of customer $i$, $\forall i \in I$}, \\
C^F_j & = & \mbox{capacity of facility $j$, $\forall j \in J$}, \\
C^V  & = & \mbox{capacity of a vehicle}, \\
F_j & = & \mbox{fixed cost of opening facility $j$, $\forall j \in J$}, \\
O_{ik} & = & \mbox{operating cost of traveling arc $(i,k)$ $\forall i,k \in V$}. 
\end{eqnarray*}
\noindent\underline{Decision Variables:}
\begin{eqnarray*}
x_{ikh} & = & \left\{ \begin{array}{ll}
                 1  & \mbox{if vehicle $h$ travels directly from location $i$ to location $k$, $\forall i \in V, k \in V, h \in H$}\\
                 0  & \mbox{otherwise},                                                  
                \end{array} 
              \right. \\
y_{ih} & = & \left\{ \begin{array}{ll}
                 1  & \mbox{if vehicle $h$ visits customer $i$, $\forall i \in I, h \in H$}\\
                 0  & \mbox{otherwise},                                                  
                \end{array} 
              \right. \\
\fvarj & = & \left\{ \begin{array}{ll}
                1  & \mbox{if facility $j$ is selected to be open, $\forall j \in J$}\\
                 0  & \mbox{otherwise}.                                                  
                \end{array}  
              \right. 
\end{eqnarray*}
A vehicle flow formulation of the LRP is as follows: 
\begin{optprog}
{\bf (VF-LRP)} Minimize & \objective{\sum_{j\in J} F_j \fvarj  +  \sum_{h \in H}\sum_{i \in V}\sum_{k \in V} O_{ik}x_{ikh}} \label{obj-eb} \\
s.t. & \sum_{h\in H}\sum_{k \in V} x_{ikh} & = & 1 &  \forall  i \in I, \label{set-part}\\
           & \sum_{k\in V} x_{ikh}-\sum_{k \in V} x_{kih} & = & 0 & \forall i \in V, h \in H, \label{enterleave}\\
           & \sum_{k\in V} x_{kih} - y_{ih} & = & 0 & \forall i \in I, h \in H, \label{cust-assgn}\\   
           &  y_{ih} - \sum_{k \in S}\sum_{l \in V\backslash S}x_{klh} & \leq & 0 &  \forall S \subseteq I, i \in S, h \in H, \label{subtour}\\   
           & \sum_{i \in I}D_i y_{ih} - C^V & \leq & 0 & \forall  h \in H, \label{VehCap}\\
           & \sum_{h \in H_j}\sum_{i \in I} D_i y_{ih} - C^F_j \fvarj & \leq & 0 & \forall j \in J, \label{FacCap}\\
           & x_{ikh} & \in & \{0,1\} & \forall i,k \in V, h \in H, \label{x-var}\\
           & y_{ih} & \in & \{0,1\} &  \forall i \in I, h \in H, \label{dri-var}\\
           & \fvarj  & \in & \{0,1\} & \forall j \in J. \label{t-var}
\end{optprog}
In (VF-LRP), the objective function (\ref{obj-eb}) seeks to minimize
the total cost, which includes the fixed cost of the selected
facilities and the operating cost of the vehicles. Vehicle fixed costs
can easily be incorporated into the model by increasing the cost of
traveling from the facility to each customer location by a fixed
amount. Constraints (\ref{set-part}) specify that exactly one vehicle
must service customer $i$. Constraints (\ref{enterleave}) require that
each vehicle should enter and leave a location the same
number of times.
Constraints (\ref{cust-assgn}) determine the assignment of customers
to vehicles. Constraints (\ref{subtour}) eliminate subtours, i.e. routes that do not
include a facility. Constraints (\ref{VehCap}) are vehicle capacity
constraints. Constraints (\ref{FacCap}) ensure that the capacity of
a facility is not exceeded by demand flows to customer locations. 
For notational convenience, we assume that variables associated with
travel between different facilities or travel between a customer and a
facility using a truck not associated with that facility are fixed to
zero. Constraints (\ref{x-var}), (\ref{dri-var}), and (\ref{t-var}) are
the set of binary restrictions on the variables.
\subsection{Set-Partitioning Formulation}
Here we utilize a modified version of the set-partitioning formulation for the
LRSP presented by \citet{Akca}. 
We define the set $P$ to be a set indexing all vehicle routes feasible
with respect to vehicle capacity and originating from and returning to
the same facility. We let $P_j \subseteq P$
index routes associated with vehicles assigned to facility $j$ for $j
\in J$. In addition to the parameters and sets defined for
(VF-LRP), we use the following parameters and decision variables:

\noindent\underline{Parameters:}
\begin{eqnarray*}
O_p & = & \mbox{operating cost of route $p \in P$}, \\
a_{ip} & = & \left \{  \begin{array}{ll}
                 1 & \mbox{if customer $i \in I$ is assigned to route
		   $p \in P$, and }\\ 
                 0 & \mbox{otherwise}.
                         \end{array}
              \right. 
\end{eqnarray*}
\noindent\underline{Decision Variables:}
\begin{eqnarray*}
z_{p} & = & \left \{ \begin{array}{ll}
                 1  & \mbox{if route $p \in P$ is selected, and}\\
                 0  & \mbox{otherwise},                                                  
                \end{array}
       \right. \\
\fvarj &  = & \left \{ \begin{array}{ll}
                 1  & \mbox{if facility $j \in J$ is selected to be open, and}\\
                 0  & \mbox{otherwise},                                                  
                \end{array}
       \right.
\end{eqnarray*}
The formulation is as follows: 
\begin{optprog}
{\bf(SPP)} Minimize & \objective{\sum_{j \in J} F_j \fvarj + \sum_{p\in P} O_{p} z_{p} } \label{spp-obj}\\
subject to & \sum_{p \in P} a_{ip} z_{p} & = & 1 & \forall i \in I, \label{spp_con1} \\
           & \sum_{p \in P_j} \sum_{i \in I} D_i a_{ip}  z_{p} - C^F_{j} \fvarj & \leq & 0 & \forall j \in J, \label{spp_con2} \\
           & z_{p} &  \in & \{0,1\} & \forall p \in P, \label{binary1} \\
           & \fvarj  & \in & \{0,1\} & \forall  j \in J. \label{binary2} 
\end{optprog} 
The objective function (\ref{spp-obj}) seeks to minimize the total
cost, which includes the fixed cost of the selected facilities and the
operating cost of the vehicles. Constraints (\ref{spp_con1}) guarantee
that each customer location is served by exactly one route.
Constraints (\ref{spp_con2}) ensure that the total demand of the
selected routes for a facility does not exceed the facility capacity.
Constraints (\ref{binary1}) and (\ref{binary2}) are standard
binary restrictions on the variables.

\subsection{Comparing the Formulations}
Observe that the three-index formulation (VF-LRP) exhibits a high
degree of symmetry under the assumption that the vehicle fleet
assigned to each facility is homogeneous. This is due to the fact that
the assignment of routes to a specific vehicle is essentially arbitrary,
i.e., the cost of a given solution to (VF-LRP) is invariant under
permutation of the indices assigned to specific vehicles. This
symmetry can be dealt with either (i) by using a two-index
formulation, which requires the addition of an exponential number of
valid inequalities to the formulation or (ii) by applying
Dantzig-Wolfe decomposition (DWD). \citet{Laporte86} and
\citet{Belenguer} have each developed branch-and-cut algorithms using
the former approach. Here, we explore the latter approach.

As is standard in DWD, we decompose the constraints of (VF-LRP) into
two subsystems. The \emph{master problem} is defined by constraints
(\ref{set-part}) and (\ref{FacCap}), while the \emph{subproblem} is
defined by constraints (\ref{enterleave}), (\ref{cust-assgn}),
(\ref{subtour}), (\ref{VehCap}), (\ref{x-var}), (\ref{dri-var}), and
(\ref{t-var}). All constraints of the subproblem are indexed by the
set $J$ and it is therefore immediate that the subproblem decomposes
by facility. The objective of each of the resulting single-facility
subproblems is to generate least-cost routes for all vehicles assigned
to the facility, though without the constraint that customers appear
on exactly one route. For each candidate facility $j \in J$, the set
of integer solutions of the decomposed subproblem can be represented
by the set of vectors
\begin{eqnarray*}
\{(x, y, t) \in \Z^{(|V|\times |V| \times |H_j|) \times (|V| \times
  |H_j|) \times \{0, 1\}}  \mid & &   (x, y, t) \mbox{ satisfies } (\ref{enterleave})_j, 
 (\ref{cust-assgn})_j, (\ref{subtour})_j, \\
& & (\ref{VehCap})_j, (\ref{x-var})_j,
(\ref{dri-var})_j, (\ref{t-var})_j\},
\end{eqnarray*}
which is described only by those constraints associated with facility
$j$ and specifies values only for those variables associated with
vehicles assigned to facility $j$. Hence, the index $j$ for the
constraints represents the set of constraints associated with facility
$j$. Let $E$ be a set indexing the members of all of the above sets,
with $E_j$ the indices for vectors associated with facility $j$ only,
so that $E = \cup_{j \in J} E_j$. For a facility $j$ and an index $q
\in E_j$, the corresponding member $(\mbox{x}^q, \mbox{y}^q,
\mbox{t}^q)$ of the above set is then a vector with the following
interpretation: for each pair $(i, k) \in |V| \times |V|$ and $h \in
H_j$, the parameter $\mbox{x}_{ikh}^q = 1$ if vehicle $h$ travels on
arc $(i,k)$ in solution $q$ and is 0 otherwise; for each $i \in I$ and
$h \in H_j$, $\mbox{y}_{ih}^q = 1$ if customer $i$ is visited by
vehicle $h$ in solution $q$ and is 0 otherwise; and $\mbox{t}^q =1$ if
facility $j$ is open in solution $q$. Note that the variable $t$
indicating whether the facility is open does not appear in any of the
linear constraints of the subproblem and can hence be set to either 0
or 1 without affecting feasibility. 

Because the subproblem decomposes as described above, solutions to the
original problem can be seen as vectors obtained by ``recomposing'' convex
combinations of the members of $E_j$ for each $j \in J$. In other words, any
solution $(x, y, t)$ to the LP relaxation of the original problem can be
written as:
\begin{optprog}
& x_{ikh} & = & \sum_{q \in E}\mbox{x}_{ikh}^q \theta_q &  \forall  i,k \in V, h \in H,  \label{XintermsTeta}\\
& y_{ih} & = &  \sum_{q \in E}\mbox{y}_{ih}^q \theta_q & \forall  i \in I, h \in H, \label{TotalYintermsTeta}\\
& t_j  & = & \sum_{q \in E_j}\mbox{t}^q\theta_q   &  \forall  j \in J, \label{Tvartrans}\\
& \sum_{q \in E_j}\theta_q & = & 1 & \forall  j \in J, \label{eq1}\\
& \theta_q  & \geq & 0 &  \forall q \in E. \label {conv}
\end{optprog}  
Using (\ref{XintermsTeta}) - (\ref {conv}), we can then formulate the
LP relaxation of the master problem as:
\begin{optprog}
{\bf (MDW)} Minimize & \objective{\sum_{q \in E} \tilde{C_{q}} \theta_q} \nonumber\\
s.t. & \sum_{q \in E} b_{iq} \theta_q & = & 1 & \forall i \in I, \label{spar} \\
           & \sum_{q \in E_j} \sum_{i \in I} b_{iq} D_i \theta_q  -
C^F_j \sum_ {q \in E_j}\theta_q \mbox{t}^q & \leq & 0 & \forall j \in J, \label{rel-2}  \\
           & \sum_{q \in E_j} \theta_q  & = & 1 & \forall j \in J,\label{conv2} \\
           & \theta_q  & \geq & 0 &  \forall q \in E, \label {conv3}
\end{optprog} 
where 
\begin{eqnarray*}
 b_{iq} & = & \sum_{h \in H}\sum_{k \in I}\mbox{x}_{ikh}^q = \sum_{h \in H}\mbox{y}_{ih}^q, \hspace{0.2in} \forall i \in I, q \in E, \label{transform-xa}\\
\tilde{C_{q}} & = & F_q {\mbox{t}}^q+ \sum_{i \in V} \sum_{k \in V}\sum_{h \in H} O_{ik}\mbox{x}_{ikh}^q, \hspace{0.2in} \forall  q \in E, \label{transform-cost} \\
\end{eqnarray*}
where $F_q = F_j$ when $q \in E_j$ for $j \in J$. Here, $ b_{iq}$ can be
interpreted as the number of times customer $i$ is visited in solution
$q$, and $\tilde{C_{q}}$ is the cost of solution $q$ (including
facility fixed cost) for all $q\in E$. 

The similar forms of (SPP) and (MDW) should now be evident, but to
rigorously show their equivalence, we need to dissect the relationship
between set $E_j$ and $P_j$ for a given facility $j \in J$. A member
of set $E_j$ consists of a collection of routes assigned to vehicles
located at facility $j$. A member of set $P_j$, on the other hand, is
a single route that can be assigned to any vehicle at facility $j$.
Therefore, a member of $E_j$ can be constructed by associating at most
$|H_j|$ members of set $P_{j}$ and some number of empty routes (zero
vectors representing vehicles that are not used) to the vehicles
assigned to facility $j$.

Now, by utilizing the integrality requirements from the original
problem and carefully eliminating the indices of symmetric solutions
from $E$, we get a much smaller set that we will show is in one-to-one
correspondence with collections of members of $P_j$ of cardinality at
most $|H_j|$. We proceed as follows:

\begin{enumerate}
\item 
First, as the vehicles associated with a given facility $j$ are
identical, a set of routes from $P_j$ can be assigned to vehicles in
any arbitrary order and each route can also visit the customers in
either clockwise or counterclockwise order. Hence, we obtain
different members of $E_j$ that are all equivalent from the standpoint
of both feasibility and cost.
To eliminate superfluous equivalent members of $E_j$, we divide the
members of set $E_j$ into equivalence classes, where two members of
$E_j$ are considered equivalent if the set of customers assigned to
the facility and the partition of that set of customers defined by the
routes are identical. It is clear that any two members of $E_j$ that
are equivalent by this definition will have exactly the same impact on
both cost and feasibility. We then form equivalence classes from which
all but one member may safely be eliminated from $E_j$.

\item Let $\hat{\theta}$ represent a solution to (MDW) for which the
  corresponding solution in the original space obtained by applying
  equations (\ref{XintermsTeta})-(\ref{Tvartrans}) is feasible for the
  original problem (VF-LRP). From (\ref{XintermsTeta}) and
  (\ref{x-var}), along with the fact that $\mbox{x}^q_{ikh}$ is binary
  for $q \in E$, we get that $\mbox{x}^q_{ikh} = 1$ whenever
  $\hat{\theta}_q > 0$ in the solution to (MDW). In other words, all
  members $q$ of $E_j$ with $\hat{\theta}_q > 0$ must correspond to
  routes visiting exactly the same customers in exactly the same
  order. Hence, we must in fact have $\hat{\theta}_q = 1$ for exactly
  one $q \in E_j$ for each $j \in J$ and so $\theta_q \in \{0, 1\}$
  for all $q \in E$.
\item It then follows easily that index $q$ can be removed from $E_j$
if there exist vehicles $h_1, h_2 \in H_j$ such that
$\mbox{x}_{ikh_1}^q = \mbox{x}_{ikh_2}^q$ $\forall i, k \in V$, where
$\mbox{x}_{ikh_1}^q>0$ for some $i, k \in I$ (i.e., this is not an
empty route). In this case, vehicles $h_1$ and $h_2$ define exactly
the same set of routes, which means that $b_{iq}> 1$ for some $i \in
I$. Because of constraint (\ref{spar}), such a solution must have
$\hat{\theta}_q=0$.
\item Finally, from (\ref{Tvartrans}) and (\ref{rel-2}), we can
  conclude that if $\hat{\theta}_q =1$ and we have $\mbox{x}_{ikh}^q>0$
  for some $i, k \in V$ and $h \in H$, then $\mbox{t}^q$ must be $1$.
  Hence, we can eliminate any solutions for which $\mbox{t}^q = 0$
  that does not correspond to a zero solution (i.e., closed facility).
  All of this allows us to rewrite (\ref{rel-2}) in the form
\begin{equation}
\sum_{q \in E_j} \sum_{i \in I} b_{iq} D_i \theta_q  -
C^F_j \sum_{q \in E_j} \theta_q \leq 0 \ \ \forall j \in J. \label{rel-2-new}
\end{equation}
\end{enumerate}   
If we restrict set $E_j$ according to the rules described above and call the
restricted set $\bar{E}_j$ for each $j \in J$, then we can finally
conclude the following.
\begin{proposition}
\label{dw-prop-p1}
There is a one-to-one correspondence between subsets of $P_j$ with
cardinality less than $|H_j|$ and members of $\bar{E}_j$.
\end{proposition}
The proof follows easily from the definition of sets $P_j$ and $\bar{E}_j$
and the restriction rules. By replacing set $E_j$ with $\bar{E}_j$
for all $j \in J$ in (MDW), as well as replacing (\ref{rel-2}) with
(\ref{rel-2-new}) and adding the constraint $\theta_q \in \{0, 1\}$ for
all $q \in E$, we obtain a new (equivalent) formulation (MDW'). Then,
we finally have the equivalence of (SPP) and (MDW) as follows
\begin{proposition}
\label{dw-prop-p2}
There is a one-to-one correspondence between solutions to (SPP)
and solutions to (MDW') such that corresponding solutions also have the
same objective function value. Thus, (SPP), (MDW'), and (MDW) are all
equivalent.
\end{proposition}

\section{Solution Algorithm}
\label{sec:SolAlg}
\subsection{Branch-and-Price}
Having shown that (SPP) is equivalent to a DWD
of (VF-LRP), we now discuss an exact solution algorithm based on a
branch-and-price implementation utilizing the formulation (SPP).
First, we strengthen the original formulation by adding the
following additional valid constraints:
\begin{optprog}
           & \sum_{p \in P_{j}}a_{ip}z_{p} - \fvarj &\leq & 0 & \forall i \in I, \forall j \in J, \label{relation}\\
           & \sum_{j \in J} \fvarj & \geq & N^F, \label{lb-open-fac} \\
           & \sum_{p \in P_j}z_p & = & v_j  &\forall j \in J, \label{veh_var} \\
           & \sum_{j \in J} v_j & \geq & \left \lceil \frac{\sum_{i \in I} D_i}{C^V} \right \rceil \label{vec_bnd} \\
           & v_j & \geq & t_j  &\forall j \in J, \label{vec_fac} \\
           & v_j  & \in & \mathbb{Z}^+ & \forall  j \in J, \label{vint}  
\end{optprog} 
where $v_j$ represents the number of vehicles used at facility $j$,
$\forall j \in J$ and $N^F$ is a lower bound on the number of facilities
that must be opened in any integer solution, calculated as follows:
\begin{eqnarray*}
N^F  & = & \mbox{argmin}_{\{l = 1,...,|J|\}} \left (\sum_{t=1}^{l}C^F_{j_t}\geq \sum_{i \in I} D_i \right ) \ \ \mbox{s.t. }  C^F_{j_1}\geq C^F_{j_2} \geq ...\geq C^F_{j_n}.
\end{eqnarray*}
Constraints (\ref{relation}) force a facility to be open if any
customer is assigned to it. Constraint (\ref{lb-open-fac}) sets a
lower bound on the total number of facilities required in any integer
feasible solution.
Constraints (\ref{relation}) (from \citet{Berger07} and \citet{Akca})
and constraint (\ref{lb-open-fac}) (from \citet{Akca}) are shown
computationally to improve the LP relaxation of the model.
Constraints (\ref{veh_var}) are only added to facilitate branching
on the integrality of the number of vehicles at each facility in the
branch-and-price algorithm. Constraint (\ref{vec_bnd}) sets a lower
bound for the total number of vehicles in the solution.
Finally, constraints (\ref{vec_fac}) force the number of vehicles used
at a facility to be at least 1 if the facility is open. We refer to
formulation (\ref{spp-obj})--(\ref{binary2}) and
(\ref{relation})--(\ref{vint}) as (SP-LRP) in the rest of the paper.

The formulation (SP-LRP) contains a variable for each possible vehicle
route originating from each facility. Hence, the number of routes will
be too large to enumerate even for small instances. To solve the LP
relaxation of models that contain exponentially many columns, column
generation algorithms can be used. By solving a sequence of linear
programs and dynamically generating columns eligible to enter the
basis, such an algorithm implicitly considers all columns but generates
only those that may improve the objective function. In order to
generate integer solutions, a branch-and-bound approach is used in
combination with the column generation and the overall approach is
referred to as branch-and-price. \citet{Desrochers2}, \citet{Vance2},
\citet{Berger07}, and \citet{Akca} provide examples of
branch-and-price algorithms from the literature.

Here, we modify the branch-and-price algorithm used in \citet{Akca} to
solve the LRSP. Therefore, some parts of the algorithm are described
only briefly. For details, we refer the reader to \citet{Akca}. To
initialize the algorithm, we construct a \emph{restricted master
problem} (RMP), that is, an LP relaxation of (SP-LRP) that contains
all facility variables ($t_j$ for $j \in J$) and vehicle variables
($v_j$ for $j \in J$), but only a subset of the vehicle route
variables ($z_p$ for $p \in P$). The branch-and-price
algorithm consists of two main components: an algorithm for
solving the \emph{pricing problem} or \emph{column generation
subproblem}, which is used to construct new columns in each iteration,
and \emph{branching rules}, which specify how to partition the
feasible region into subsets to which the algorithm is then applied
recursively until exhaustion.

At each iteration of the solution process for the LP relaxation of
(SP-LRP), the objective of the column generation subproblem is to find
a feasible vehicle route originating from each facility $j \in J$ with
minimum reduced cost with respect to the current dual solution of the
RMP. The reduced cost of a given route $p \in P$ is
\begin{eqnarray}
\label{red-cost}
\hat{C}_{p}=O_{p}-\sum_{i\in N}a_{ip}\pi_{i}+\sum_{i \in N}a_{ip}
D_{i} \mu_{j}+\sum_{i \in N}a_{ip} \sigma_{ji} - \nu_j,
\end{eqnarray}
where $\pi$, $\mu$, $\sigma$ and $\nu$ are the dual variables
associated with constraints (\ref{spp_con1}), (\ref{spp_con2}),
(\ref{relation}), and (\ref{veh_var}), respectively, from the RMP.
Hence, the column generation subproblem for facility $j \in J$ can be
formulated as follows:
\begin{optprog}
Minimize & \objective{\sum_{i \in M}(\sum_{k \in I} (O_{ik}-\pi_{k}+D_k \mu_j + \sigma_{jk})x_{ik} + O_{ij}x_{ij})- \nu_j}  \label{obj-pri} \\
s.t.       & y_j & = & 1, \\ 
           & \sum_{k \in M} x_{ik} = \sum_{k \in M} x_{ki} & = & y_i &  \forall  i \in I, \label{set-part-pri}\\
           &  y_{i} - \sum_{k \in S}\sum_{l \in V \backslash S}x_{kl} & \leq & 0 &  \forall S \subseteq I, i \in S, \label{subtour-pri}\\   
           & \sum_{i \in I}D_i y_{i} - C^V & \leq & 0, \label{VehCap-pri}\\
           & x_{ik} & \in & \{0,1\} & \forall i,k \in M, \label{x-var-pri}\\
           & y_{i} & \in & \{0,1\} &  \forall i \in M, \label{dri-var-pri}
\end{optprog}
where $M = I \cup \{j\}$, $x_{ik} = 1$ if link $(i,k)$ is used in
solutions, and $y_i = 1$ if customer $i$ is assigned to the route.

The column generation subproblem above can be seen as an instance of the
well-known \emph{elementary shortest path problem with resource
constraints} (ESPPRC), which is well-studied and arises as a
subproblem in many different routing applications. To cast the above
subproblem as an ESPPRC, we consider a network with $|I| + 2$ nodes,
one node for each customer, one for the facility $j$ as a source node
and a copy of the facility $j$ as a sink node. We assign each customer
node a demand equal to its demand in the original problem and let the
cost of each arc $(i,l)$ in the network equal the coefficient of
$x_{il}$ in (\ref{obj-pri}). A shortest path from source to sink
visiting a customer at most once (called \emph{elementary}) and
satisfying the constraint that the total demand of customers included
in the path does not exceed the vehicle capacity then corresponds to a
vehicle route. The total cost of the path plus the fixed cost $-\nu_j$
is the reduced cost of the associated column. 

To solve the ESPPRC, we use \citet{Feillet}'s label setting algorithm
with a single resource (the vehicle capacity). In the algorithm, each
path from source to sink that is not dominated by another with respect
to vehicle capacity, cost, and the set of nodes that can still be
visited is explored. More details on the variants of this approach
that were used in the computational experiments are given in
Section~\ref{sec:pricing} below.

When the pricing problem cannot identify any more columns with
negative reduced cost, then the current solution to the LP relaxation
of the master problem is optimal. If this optimal LP solution is not
integral, then we are forced to branch. Devising good branching rules
is an important step in developing a branch-and-price algorithm. Since
the LP relaxations of the new nodes generated after branching are also
solved using column generation, the branching constraints must be
incorporated into the pricing problem and the columns to be generated
must satisfy branching constraints for the node. Therefore, the
specific branching rules employed may affect the structure of the
pricing problem, causing it to become more difficult to solve. Here,
we implement the same four branching rules we used for our work on the
LRSP: branching on fractional variables $t$ and $v$, forcing/forbidding the
assignment of a specific customer to a specific facility, and
forcing/forbidding flow on a single arc (originally used in
\citet{Desrochers4}). All branching rules
can easily be incorporated into the pricing problem without
changing the structure. Details on the effect of these rules on the
pricing problem and the implementation of them can be found in
\citet{Akca}.

\subsection{Solving the Column Generation Subproblem}
\label{sec:pricing}
The ESPPRC is an NP-hard optimization problem, but small
instances may generally be solved effectively in practice using
dynamic programming-based labeling algorithms. In general, the number
of labels that must be generated and evaluated in the label setting
algorithm increases as either the number of customers or the vehicle
capacity increases. To enhance efficiency, we therefore augment the
basic scheme with some heuristic versions of the algorithm, since it
is not necessary to find the column with the most negative reduced
cost in every iteration. We refer to the exact pricing algorithm (that is
guaranteed to produce a column with the smallest reduced cost) as ESPPRC.
The following are heuristic versions of the exact algorithm and are
not guaranteed to produce a column with negative reduced cost when one
exists.

%\paragraph{ESPPRC-LL(n).} 
{\bf ESPPRC-LL(n).} The ESPPRC algorithm keeps all non-dominated labels
at each node. Depending on the size of the instance, the number of labels 
kept can become very large. 
In this heuristic version proposed by \citet{Dumitrescu}, we set a limit $n$ on the number of
unprocessed labels stored at each node. At each iteration, labels are
sorted based on reduced cost and among the unprocessed labels, the
$n$ with smallest reduced cost are kept and the rest are permanently
deleted. Therefore, the ESPPRC-LL(n) algorithm tends to terminate much
more quickly than the ESPPRC for small values of $n$.

%\paragraph{ESPPRC-CS(n).} 
{\bf ESPPRC-CS(n).} The ESPPRC algorithm is efficient for instances
with small numbers of customers. In addition, the total number of
customers in a route is restricted by the vehicle capacity. To be able
to take advantage of this, we choose a subset of customers $C_S$ with
size $n$ based on the cost of arcs in the network 
(coefficients of link variables in (\ref{obj-pri})).  Since the
objective of the pricing problem is to find a route with smallest
reduced cost, we determine the $n$ customers in the subset by
constructing a minimum spanning tree over the customer locations. We stop
the spanning tree algorithm when we have $n$ customers in the tree.
The first customer in the subset (and the tree) is chosen based on the
cost of the links from source to customers. Then to find valid vehicle
routes, we run ESPPRC that include only customers in $C_S$.

%\paragraph{2-Cyc-SPPRC-PE(n).}  
{\bf 2-Cyc-SPPRC-PE(n).} The shortest path problem with resource
constraints (SPPRC) is a relaxation of the ESPPRC in which the path
may visit some customers more than once.
The SPPRC is solvable in pseudo-polynomial time, but use of this
further relaxation of the column generation subproblem results in
reduced bounds. Eliminating solutions containing cycles of length two
strengthens this relaxation of the pricing problem and improves the
bound (for details of the algorithm, see
\citet{IrnichDesaulniers2005}). We refer to this
pricing algorithm as 2-Cyc-SPPRC. In addition, we can also generate
paths that are elementary with respect to a given subset of customers.
We call the resulting algorithm 2-Cyc-SPPRC-PE(n), where $n$ is the
size of the customer set to be considered.
At each iteration of the pricing problem, for each
facility, the algorithm consists of the following steps:
\begin{enumerate}
\item Solve the 2-Cyc-SPPRC.  
\item Consider the column with minimum cost and choose at most $m$
customers that are visited more than once. Let $C_{E1}$ be the set of
these customers. If the set is empty, stop (the path is already
elementary).
\item Solve 2-Cyc-SPPRC-PE(m) with set $C_{E1}$ from step 2. 
\item Pick at most $m$ customers that are visited more than once. Let $C_{E2}$
be the set of these customers. Let $C_{E3} = C_{E1}\cup C_{E2}$ and $m_3 =
|C_{E3}|$. If the set is empty, let $C_{E3} = C_{E1}$ and stop.
\item Solve  2-Cyc-SPPRC-PE($m_3$) with set $C_{E3}$.
\end{enumerate} 
In our experiments, we generally had the same set of customers
$C_{E3}$ in every step of the column generation. Thus, we decided to
determine set $C_{E3}$ for each facility at the first iteration of
column generation at each node, and we use the same set for the rest
of the iterations at the node.
  
For computational testing, we implemented four variants of the
branch-and-price algorithm based on the above pricing schemes.
\begin{itemize}
\item \emph{Heuristic Branch-and-Price} (HBP): The purpose of this
algorithm is to provide a good upper bound. At each node of the tree,
we use ESPPRC-CS(n) and ESPPRC-LL(m) with small values of $n$ and
$m$ ($n$ is chosen to be 12 to 15 depending on the average demand and
the vehicle capacity, while $m$ is chosen to be 5 or 10). In addition, we
also used combinations of ESPPRC-CS(n) and ESPPRC-LL(m) for
larger values of $n$. In the algorithm, we use an iteration limit for
the number of pricing problems solved at any node of the tree. If the
number of iterations exceeds the limit, we branch.

\item \emph{Elementary Exact Branch-and-Price} (EEBP): The purpose of
this algorithm is to prove the optimality of the solution or provide
an integrality gap. At each node of the tree,
we use ESPPRC-CS(n), ESPPRC-LL(m) and ESPPRC.

\item \emph{2 Step Elementary Exact Branch-and-Price} (EEBP-2S): In
  this variant, HBP is run first to generate initial columns and an
  upper bound. Then, EEBP is initiated with the columns and the upper
  bound obtained from HBP.

\item \emph{Non-elementary Exact Branch-and-Price (NEBP)}: This is
similar to elementary exact branch-and-price algorithm except that
the pricing problem solved is 2-Cyc-SPPRC-PE(n).

\end{itemize}

\section{Computational Results}
\label{sec:CompR}
In this section, we discuss the performance of our branch-and-price
algorithm for the LRP on two sets of instances. The first set contains the
LRP instances used in \citet{Barreto07} that are available from
\citet{Barreto-data}. 
We used these instances to test the performance of our HBP and NEBP
algorithms. The second set of instances were random instances we
generated to test the performance of our EEBP and EEBP-2S algorithms.
We evaluate the effect of facility capacity constraints and other
parameters using this set of instances. For all of our experiments, we
used a Linux-based workstation with a 1.8 GHz processor and 2GB RAM.

\subsection{Instances From the Literature}
To the best of our knowledge, there are no benchmark instances
available specifically for the LRP. \citet{Barreto07} used the
instances in the literature available for other types of problems to
construct a set of LRP instances. They report lower bounds found by
applying a branch-and-cut algorithm to the two-index formulation of
the problem \citep{Barreto04} and upper bounds found by applying
a sequential heuristic based on clustering techniques
\citep{Barreto07}. They listed 19 instances, three of which have
more than 150 customers, too large for our approach to work
efficiently. We removed these three instances plus one more with 117
customers and fractional demand, since we assume integer demands. The
labels of the instances denote the source of the instance and the
number of customers and facilities in the instances
(for more details about the references, see \citet{Barreto-data}).

We first ran HBP with a time limit of 3 CPU hours, focusing on
producing quality upper bounds. 
Table \ref{brt-ub} presents the instances we tested and compares the
results with the upper bounds reported in \citet{Barreto07}. Since
neither our HBP nor \citet{Barreto07} can provide a valid lower bound
for the problem, we used the best lower bounds as found in
\citet{Barreto04} (second column in Table \ref{brt-ub}) to measure the
quality of our upper bound. The ``Gap'' in Table \ref{brt-ub} is the
percent gap between the upper bound and the LB listed in the second column. HBP
is capable of finding better upper bounds (usually optimal) for the
instances of small and medium size. In these cases, the computation
time is also very short. However, for larger instances, the upper
bounds reported by \citet{Barreto07} are generally better. In
addition, their heuristic algorithm is very efficient---they report
that in most of the instances, it provides the result in less than one
second.

\begin{table}[ht]
\centering
\caption{Performance of  Heuristic Branch-and-Price}\label{brt-ub}
\begin {tabular}[t]{|l|c|c|c|c|c|c|}  
\hline
Instance & LB$^1$ & \multicolumn{2}{|c|}{\citet{Barreto07}} &  \multicolumn{3}{|c|}{Heuristic Branch-and-Price} \\\cline{3-7}
\mbox{} & \mbox{} & UB$^2$  & Gap & UB & Gap & CPU(s)\\\hline
Gaskell67-21x5	&	424.9$^*$	&	435.9	&	2.59	&	424.9	&	0	&	1.2	\\
Gaskell67-22x5 	&	585.1$^*$	&	591.5	&	1.09	&	585.1     &	0	&	41.5	\\
Gaskell67-29x5 	&	512.1$^*$	&	512.1	&	0	&	512.1	        &	0	&	67.7	\\
Gaskell67-32x5 	&	556.5	        &	571.7	&	2.73	&	562.3	&	1.04	&	10801.8	\\
Gaskell67-32x5-2 \ \ &	504.3$^*$	&	511.4	&	1.41	&       505.8	&	0.3	&	85.6	\\
Gaskell67-36x5 	&	460.4$^*$	&	470.7	&	2.24	&	460.4     &	0	&	1077.6	\\
C69-50x5 	&	549.4	        &	582.7	&	6.06	&	565.6	&	2.95	&	239.4	\\
C69-75x10 	&	744.7	        &	886.3	&	19.01	&	852.1	&	14.42	&	10802.3	\\
C69-100x10 	&	788.6	        &	889.4 &	12.78	&	929.5     	&	17.86	&	10836.6	\\
Perl83-12x2 	&	204$^*$	        &	204	&	0	&	204.0	        &	0	&	0.2	\\
Perl83-55x15 	&	1074.8 \ \	        &	1136.2	&	5.71	&	1121.8  &	4.37	&	10800.0	\\
Perl83-85x7 	&	1568.1	        &	1656.9 &	5.66	&	1668.2	        &	6.38	&	10813.8	\\
Min92-27x5 	&	3062$^*$	&	3062	&	0	&	3062.0	        &	0	&	4.2	\\
Min92-134x8 	&	  -	        &	6238 &	-	&	6421.6	        &	-	&	10850.9	\\
Dsk95-88x8 	&	356.4	        &	384.9 &	8	&	390.6	        &	9.58	&	10808.9	\\
\hline
\multicolumn{7}{l}{$^1$ Reported by \citet{Barreto04}, found using branch-and-cut}\\
\multicolumn{7}{l}{$^2$ Reported by \citet{Barreto07}, found using a heuristic}\\
\multicolumn{7}{l}{$^*$ Branch-and-cut used in \citet{Barreto04} proves the optimality}\\
\end{tabular}
\end{table}

Next, we used the EEBP-2S algorithm to test the ability to produce
lower bounds and prove optimality. With this algorithm, we could not
solve all of the instances within the total time limit of 5 CPU hours (3 CPU 
hours for HBP and 2 CPU hours for EEBP). The
lower bounds found by our algorithm, along with the best integer
solution found, optimality gap, and computation time are reported in
Table \ref{brt-gr2}. Note that because Gaskell67-21x5 and Perl83-12x2
are very easy to solve, we used the EEBP algorithm instead of the
EEBP-2S algorithm in these cases.
\begin{table}[ht]
\centering
\caption{Performance of 2 Step Elementary Exact Branch-and-Price}\label{brt-gr2}
\begin {tabular}[t]{|l|c|c|c|c|c|}  
\hline
Instance	        &       LB	&	OPT/BestIP &	Gap    & CPU(s)	& 	Total CPU(s)	\\\hline
Gaskell67-21x5$^1$	&	424.9	&	424.9	&	0	&	3.0	&	3.0	\\
Gaskell67-22x5   	&	585.1	&	585.1	&	0	&	2999.9	&	3041.4	\\
Gaskell67-32x5  	&	544.1	&	562.3	&	3.24\%	&	8453.1	&	19254.9	\\
Perl83-12x2$^1$ 	&	204.0	&	204.0	&	0	&	0.2	&	0.2	\\
Min92-27x5      	&	3062.0 \ \ &	3062.0	&	0	&	833.6	&	837.8	\\
\hline
\multicolumn{6}{l}{$^1$ EEBP algorithm is used}\\
\end{tabular}
\end{table}

Finally, for instances that we could not provide lower bounds by using
the EEBP or the EEBP-2S algorithm, we used the NEBP algorithm with a
time limit of 5 CPU hours or evaluated node limit of 50 nodes. The
results are reported in Table \ref{brt-gr3}. For the instances
C69-100x10, Min92-134x8 and Dsk95-88x8, we could not solve even the
root node within the time limit. 

In general, the lower bounds found by using branch-and-cut \citep{Barreto04} 
are better than our lower bounds. However, in some
cases, their computation times are much larger than our time limits.
The HBP algorithm can provide good upper bounds, but for the medium
and large size problems, we need to improve our lower bounding,
perhaps by adding dynamic cut generation to our algorithm in order to
close the optimality gap.
\begin{table}[ht]
\centering
\caption{Performance of Non-elementary Exact Branch-and-Price}\label{brt-gr3}
\begin {tabular}[t]{|l|c|c|c|c|}  
\hline
Instance	        &       LB	&	OPT/BestIP &	Gap    & CPU(s)\\\hline
Gaskell67-29x5 	&	441.2	&	512.1	&	13.85\%	&	14411.5	\ \ \\
Gaskell67-32x5-2 \ \	&	494.4	&	505.8	&	2.26\%	&	5654.8	\\
Gaskell67-36x5 	&	455.5	&	460.4	&	1.05\%	&	100.1\\
C69-50x5 	&	526.2	&	565.6	&	6.96\%	&	2634.1\\
C69-75x10 	&	693.5	&	852.1	&	18.62\%	&	8785.8\\
Perl83-55x15 	&	852.3	&	1121.8	&	24.02\%	&	319.5\\
Perl83-85x7 	&	1272.4	&	1668.2	&	23.73\%	&	1379.6\\
\hline
\end{tabular}
\end{table}

\subsection{Random Instances}
On random instances, \citet{Laporte86} provided computational results
for an exact method (branch-and-cut algorithm) for the capacitated
LRP.
\citet{Belenguer} also developed a branch-and-cut algorithm for the
capacitated LRP, but neither the details of the instances nor the
computational results are publicly available. Therefore, we evaluated
our algorithm by generating random instances as in \citet{Laporte86},
where they generated instances with 10, 15 and 20 customers and 4 to 8
facilities. In addition to these, we generated instances with 30 and
40 customers and 5 facilities to test the performance of the algorithm
on larger instances. The coordinates of the customers and the
facilities and the demand of each customer were generated using a
Uniform distribution on [0,100]. We then calculated the Euclidean
distance between each pair of customers and between customers and
facilities and rounded the calculated distance to two decimal places.
Demand for each customer was rounded to the nearest integer. Vehicle
capacity $C^V$ was calculated as
\begin{equation}
C^V = (1-\alpha) max_{i \in I}\{D_i\} + \alpha \sum_{i \in I}D_i,
\end{equation}   
where $\alpha$ was a parameter and the  values were chosen in set \{0,
0.25, 0.5, 0.75, 1\}.  

\subsubsection{Small and Medium Random Instances}
\citet{Laporte86} solved location and routing problems with capacitated
vehicles, but they did not have a facility capacity. Instead, they had
a lower and upper bound for the total number of facilities that could
be open in a solution. In this experiment, in order to provide a
better comparison of our algorithm with that of \citet{Laporte86}, we
removed constraint (\ref{spp_con2}) from the SP-LRP. We set $N^F$, the
minimum number of open facilities in (\ref{lb-open-fac}), to 1, and we
added constraint $\sum_{j \in J} t_{j} \leq M^F$, where $M^F$ be the
maximum number of facilities that can be open in any solution.
Facility and vehicle fixed costs were set to be zero. As in
\citet{Laporte86}, three groups of instances with different numbers of
customers and facilities were available. For each group, five
different vehicle capacities were calculated by changing $\alpha$.
Some details about the instances are listed in Table \ref{inst1}.
\begin {table}[ht]
\begin{center}
\caption {Details for the instances} \label{inst1}
\begin {tabular}[t]{|c|c|c|c|c|c|} \hline
\# of Customers & \# of Facilities & $N^F$ & $M^F$ & \# of instances & $\alpha$ \\ \hline
10 & 4 & 1 & 3 & 3 & \{0, 0.25, 0.5, 0.75, 1\} \\
15 & 6 & 1 & 4 & 3 & \{0, 0.25, 0.5, 0.75, 1\} \\
20 & 8 & 1 & 5 & 3 & \{0, 0.25, 0.5, 0.75, 1\} \\\hline
\end{tabular}
\end{center}
\end{table}

Tables \ref{gr1-10}, \ref{gr2-15} and \ref{gr3-20} present the results
achieved with our branch-and-price algorithm. Instances listed in
these tables are labeled with the number of customers, facilities and
with letters \{a,b,c,d,e\} based on the $\alpha$ value used. For
example, the instance $r10x\textrm{4-a-1}$ has 10 customers, 4 facilities and
$\alpha = 0$. The integer from 1 to 3 after the letters represents the id
of the instances within the same group. The tables present the name of each
instance, the LP solution value at the root node, the optimal solution
value, the number of evaluated nodes, and the CPU time in seconds. In
these instances, we first ran the EEBP algorithm for 5 minutes and if
the algorithm did not terminate within 5 minutes, we switched to the
EEBP-2S algorithm. Table \ref{gr1-10} presents the results for
instances with 10 customers, Table \ref{gr2-15} presents instances
with 15 customers, and Table \ref{gr3-20} presents instances with 20
customers. We marked the instances with a ``+'' sign if the EEBP-2S
algorithm was used. For problems with at least 20 customers, we needed
to use the EEBP-2S algorithm. The branch-and-price algorithm was very
successful in finding the optimal solution quickly. 
In general, our computation times were much smaller than those
reported by \citet{Laporte86}, but it is difficult to make a fair
comparison, given advances in computing technology. In most of the
instances, the LP solution found by our algorithm at the root node was
the optimal.

The instances become more difficult when the vehicle capacity
increases ($\alpha$ increases) because the number of labels
generated in the pricing problem depends directly on the vehicle
capacity. \citet{Laporte86} observed the reverse effect with regard to
their branch-and-cut algorithm. The number of cuts generated increases
and the problem gets more difficult when the vehicle capacity is
small. This is most likely due to the fact that the problem structure
becomes more like that of the traveling salesman problem (TSP) as the
capacity is increased and the TSP is much easier to solve by branch
and cut than as a capacitated routing problem.

\begin{table}[ht]
\centering
\caption{Performance of Elementary Exact Branch-and-Price for 10 customer instances}\label{gr1-10}
\begin {tabular}[t]{|l|c|c|c|c|}  
\hline
Instance & LP & OPT. & \# of Nodes & CPU(s) \\\hline
r10x4-a-1 \ \	&	472.11	&	472.11	&	1	&	0.00	\\
r10x4-a-2 \ \	&	421.44	&	421.44	&	1	&	0.00	\\
r10x4-a-3 \ \	&	548.28	&	548.28	&	1	&	0.02	\\
r10x4-b-1 \ \	&	313.01	&	313.18	&	3	&	0.04	\\
r10x4-b-2 \ \	&	297.57	&	305.27	&	19	&	0.08	\\
r10x4-b-3 \ \	&	352.66	&	354.92	&	3	&	0.03	\\
r10x4-c-1 \ \	&	257.25	&	257.25	&	1	&	0.06	\\
r10x4-c-2 \ \	&	259.76	&	259.76	&	1	&	0.04	\\
r10x4-c-3 \ \	&	296.82	&	296.82	&	1	&	0.05	\\
r10x4-d-1 \ \	&	243.42	&	257.25	&	21	&	0.52	\\
r10x4-d-2 \ \	&	250.04	&	250.04	&	1	&	0.04	\\
r10x4-d-3 \ \	&	296.82	&	296.82	&	1	&	0.04	\\
r10x4-e-1 \ \	&	226.46	&	226.46	&	1	&	0.32	\\
r10x4-e-2 \ \	&	225.82	&	225.82	&	1	&	0.17	\\
r10x4-e-3 \ \	&	272.85	&	272.85	&	1       &	0.31	\\
\hline
\end{tabular}
\end{table} 

\begin{table}[ht]
\centering
\caption{Performance of Elementary Exact Branch-and-Price for 15 customer instances}\label{gr2-15}
\begin {tabular}[t]{|l|c|c|c|c|}  
\hline
Instance & LP & OPT. & \# of Nodes & CPU(s) \\\hline
r15x6-a-1 \ \	&	435.2 \ \	&	435.2	&	1	&	0.01	\\
r15x6-a-2 \ \	&	663.32 \ \	&	663.32	&	1	&	0.01	\\
r15x6-a-3 \ \	&	411.45 \ \	&	411.45	&	1	&	0.01	\\
r15x6-b-1 \ \	&	313.46 \ \	&	313.46	&	1	&	0.31	\\
r15x6-b-2 \ \	&	414.65 \ \	&	414.65	&	1	&	0.21	\\
r15x6-b-3 \ \	&	285.01 \ \	&	285.01	&	1	&	0.12	\\
r15x6-c-1 \ \	&	313.46 \ \	&	313.46	&	1	&	3.1	\\
r15x6-c-2 \ \	&	392.75 \ \	&	392.75	&	1	&	1.98	\\
r15x6-c-3 \ \	&	279.82 \ \	&	279.82	&	1	&	5.37	\\
r15x6-d-1 \ \	&	313.36 \ \	&	313.46	&	3	&	4.92	\\
r15x6-d-2 \ \	&	378.76 \ \	&	378.76	&	1	&	9.41	\\
r15x6-d-3 \ \	&	279.82 \ \	&	279.82	&	1	&	14.13	\\
r15x6-e-1 \ \	&	305.86 \ \	&	312.18	&	5	&	9.82	\\
r15x6-e-2$^+$ \ \	&	374.86 \ \	&	374.86	&	1	&	16.62$^+$	\\
r15x6-e-3 \ \	&	274.22	 \ \ &	274.22	&	1	&	300.01	\\
\hline
\multicolumn{5}{l}{$^+$ The EEBP-2S algorithm was used, total time is reported.}\\
\end{tabular}
\end{table} 

\begin{table}[ht]
\centering
\caption{Performance of Elementary Exact Branch-and-Price for 20 customer instances}\label{gr3-20}
\begin {tabular}[t]{|l|c|c|c|c|}  
\hline
Instance & LP & OPT. & \# of Nodes & CPU(s) \\\hline
r20x8-a-1 \ \	&	639.77	&	653.11	&	57	&	0.41	\\
r20x8-a-2 \ \	&	542.23	&	551.58	&	25	&	0.98	\\
r20x8-a-3 \ \	&	760.42	&	760.42	&	1	&	0.05	\\
r20x8-b-1 \ \	&	415.3	&	417.13	&	7	&	8.35	\\
r20x8-b-2 \ \	&	383.19	&	383.19	&	1	&	6.32	\\
r20x8-b-3 \ \	&	439.18	&	447.72	&	121	&	37.32	\\
r20x8-c-1$^+$ \ \	&	398.34	&	398.34	&	1	&	26.25$^+$	\\
r20x8-c-2$^+$ \ \	&	363.86	&	363.86	&	1	&	110.87$^+$	\\
r20x8-c-3$^+$ \ \	&	402.85	&	402.85	&	1	&	28.74$^+$	\\
r20x8-d-1$^+$ \ \	&	392.26	&	392.26	&	1	&	10.04$^+$	\\
r20x8-d-2$^+$ \ \	&	359.49	&	359.49	&	1	&	200.29$^+$	\\
r20x8-d-3$^+$ \ \	&	402.85	&	402.85	&	1	&	124.24$^+$	\\
r20x8-e-1$^+$ \ \	&	392.28	&	392.28	&	1	&	12.39$^+$	\\
r20x8-e-2$^+$ \ \	&	355.39	&	355.39	&	1	&	82.77$^+$	\\
r20x8-e-3$^+$ \ \	&	402.46	&	402.46	&	1	&	103.27$^+$	\\
\hline
\multicolumn{5}{l}{$^+$ The EEBP-2S algorithm is used, total time is reported.}\\
\end{tabular}
\end{table} 

To strictly differentiate the instances from those with a single
depot, we experimented with changing the value of parameter $N^F$, the
minimum number of open facilities, to 2 and ran the r15x6 and r20x8 b
and c instances. There were no significant changes in the
computational times or the number of evaluated nodes. We then added
facility capacities to the same set of problems and ran our algorithm
again. Table \ref{gr4-fc} presents the computational results, as well
as the facility capacity values used for the facilities. The capacity
value was chosen in order to require at least two open facilities. The
computational results do not show any significant difference from
those of the uncapacitated instances. For larger instances, we expect
that the LP solution times will tend to increase if the master problem
has facility capacity constraints. Adding a facility capacity to a
two-index vehicle flow formulation requires an additional set of
constraints \citep{Belenguer}, the size of which can be large.
In a branch-and-cut algorithm, it may require additional time to
generate this set of constraints.
\begin{table}[ht]
\centering
\caption{Instances with Capacitated Facilities}\label{gr4-fc}
\begin {tabular}[t]{|l|c|c|c|c|c|}  
\hline
Instance & Fac. Cap & LP & OPT. & \# of Nodes & Total Time \\\hline
cr15x6-c-1 \ \	&	600	&	299.07	&	299.07	&	1	&	1.15	\\
cr15x6-c-2 \ \	&	600	&	392.75	&	392.75	&	1	&	2.53	\\
cr15x6-c-3 \ \	&	600	&	279.82	&	279.82	&	1	&	1.86	\\
cr15x6-d-1 \ \	&	700	&	299.07	&	299.07	&	1	&	0.57	\\
cr15x6-d-2 \ \	&	700	&	378.76	&	378.76	&	1	&	6.42	\\
cr15x6-d-3 \ \	&	700	&	279.82	&	279.82	&	1	&	11.01	\\
cr20x8-c-1 \ \	&	750	&	398.34	&	398.34	&	1	&	40.46	\\
cr20x8-c-2 \ \	&	750	&	363.86	&	363.86	&	1	&	88.12	\\
cr20x8-c-3 \ \	&	750	&	402.85	&	402.85	&	1	&	15.25	\\
cr20x8-d-1 \ \	&	900	&	392.29	&	392.29	&	1	&	29.1	\\
cr20x8-d-2 \ \	&	750	&	359.49	&	359.49	&	1	&	257.94	\\
cr20x8-d-3 \ \	&	900	&	402.85	&	402.85	&	1	&	31.3	\\\hline
\end{tabular}
\end{table} 

\citet{Laporte86} report that adding facility fixed costs to the
problem makes the problem easier. We added facility fixed costs to
r15x6 and r20x8 b and c instances. The performance of the
branch-and-price algorithm was not affected in instances with 15
customers. However, for some of the 20 customer instances, the
computational times exceeded 2 CPU hours (for the results, see
\citet{Akca2}).


\subsubsection{Large Random Instances}
In this section, we present the results of applying our algorithm to
larger capacitated random instances. We generated 6 instances with 30
customers and 6 instances with 40 customers. Each instance had 5
facilities with capacity constraints. The facilities had a fixed cost
of 100. The characteristics of each instance are listed in Table
\ref{soln-tb}. The first column includes the name of each instance,
labeled based on the number of customers and facilities and the
vehicle capacity. For instances in group ``a'' the vehicle capacity
value (listed in the fourth column) was chosen to be 7 times the
average demand and for group ``b'', the vehicle capacity value was 5.5
times the average demand. Facility capacity (listed in the second
column) was chosen based on total demand such that at least two
facilities (the minimum number of facilities is listed in the third
column) should be open in an integer solution.

\begin{table}[ht]
\centering
\caption{Characteristics of the Instances and the Optimal Solutions}\label{soln-tb}
\begin {tabular}[t]{|l|c|c|c|c|c|c|c|}  
\hline
 \multicolumn{4}{|c|}{Instance} & \multicolumn{4}{|c|}{OPT/BestIP Solution Info} \\\hline
Name & Fac.& $N^F$ & $L^V$ &  \# of & \# of & Avg. \# of & \# of Cust. \\
\mbox{} & \mbox{Cap.} & \mbox{} & \mbox{} & Fac. & Vec. & Cust/route & longest route \\\hline
cr30x5a-1  \ \	&	1000	&	2	&	350	&	2	&	3,2	&	6	&	8	\\
cr30x5a-2  \ \	&	1000	&	2	&	350	&	2	&	3,2	&	6	&	7	\\
cr30x5a-3  \ \	&	1000	&	2	&	350	&	2	&	3,3	&	5	&	7	\\
cr30x5b-1  \ \	&	1000	&	2	&	275	&	2	&	3,2	&	6	&	8	\\
cr30x5b-2  \ \	&	1000	&	2	&	275	&	2	&	4,3	&	4.29	&	7	\\
cr30x5b-3  \ \	&	1000	&	2	&	275	&	2	&	3,4	&	4.29	&	6	\\
cr40x5a-1  \ \	&	1750	&	2	&	340	&	2	&	3,3	&	6.67	&	8	\\
cr40x5a-2  \ \	&	1750	&	2	&	390	&	2	&	3,4	&	5.71	&	8	\\
cr40x5a-3  \ \	&	1750	&	2	&	370	&	2	&	3,3	&	6.67	&	7	\\
cr40x5b-1  \ \	&	1750	&	2	&	275	&	2	&	3,5	&	5	&	7	\\
cr40x5b-2 \ \ 	&	1750	&	2	&	275	&	2	&	3,5	&	5	&	7	\\
cr40x5b-3  \ \	&	1750	&	2	&	325	&	2	&	5,3	&	5	&	7	\\
\hline
\end{tabular}
\end{table} 
We used the EEBP-2S algorithm in which the HBP and EEBP algorithms
were both used with a time limit of 3 CPU hours. Table \ref{gr5-fc}
presents the results of both steps. The algorithm was very successful
in finding optimal or near-optimal solutions for these larger
instances. 
Some details, such as the number of open facilities, the number of
vehicles used at each open facility, the average number of customers
in each route (vehicle), and the number of customers in the longest
route, are presented in Table
\ref{soln-tb}.
\begin{table}[ht]
%\centering
\caption{Performance of the EEBP-2S for Instances up to 40 Customers with Capacitated Facilities and Facility Fixed Cost}\label{gr5-fc}
\begin {tabular}[t]{|l|c|c|c|c|c|c|c|c|}  
\hline
Instance & \multicolumn{2}{|c|}{HBP} & \multicolumn{5}{|c|}{EEBP}& Total \\\cline{2-8}
\mbox{} & IP & CPU(s) & LP & IP & Gap & \# of N. & CPU(s) & CPU (s) \\\hline
cr30x5a-1  \ \ &	819.53	&	43.5	&	810.29	&	819.5	&	0	&	33	&	993.3	&	1036.8	     \\
cr30x5a-2 \ \ &	823.49	&	7202.3	&	790.49	&	823.49	&	2.55	&	500	&	10806.5 \ \	& 18008.8   \ \ \\
cr30x5a-3 \ \ &	702.29	&	44.2	&	687.72	&	702.19	&	0	&	51	&	917.9	&	962.1	\\
cr30x5b-1 \ \ &	880.02	&	164.3	&	865.47	&	880.01	&	0	&	251	&	6420.6	&	6585	\\
cr30x5b-2  \ \ &	825.32	&	8.3	&	815.95	&	825.3	&	0	&	7	&	33.2	&	41.5	\\
cr30x5b-3 \ \ &	884.62	&	13.4	&	881.33	&	884.55	&	0	&	19	&	41.7	&	55.1	\\
cr40x5a-1 \ \ &	928.11	&	631.8	&	911.39	&	928.11	&	1.49	&	11	&	10882.8 &	11514.6	\\
cr40x5a-2 \ \ &	888.37	&	378.4	&	871.66	&	888.37	&	0.93	&	13	&	11052.9 &	11431.3	\\
cr40x5a-3 \ \ &	947.24	&	173	&	939.54	&	947.24	&	0.18	&	28	&	10862 &	11035	\\
cr40x5b-1 \ \ &	1052.07 \ \	&	257.3	&	1043.62	&	1052	&	0	&	627	&	8084.6	&	8342	\\
cr40x5b-2 \ \ &	981.52	&	60.1	&	976.88	&	981.27	&	0	&	47	&	862.5	&	922.7	\\
cr40x5b-3 \ \ &	964.32	&	62.6	&	959.05	&	964.23	&	0	&	45	&	963	&	1025.6	\\
\hline
\end{tabular}
\end{table} 

\section{Conclusion}
\label{sec:Conc}
We have presented a set-partitioning-based formulation for the
capacitated location and routing problem, which to our knowledge is
the first of its kind for this class of problem. We have demonstrated
that it can be obtained by applying Dantzig-Wolfe decomposition to the
graph-based formulation employed in most previously reported research.
We have described a branch-and-price algorithm and reported on our
experience using it to solve both problems from the literature and
randomly generated instances. Our experiments indicate that the
algorithm is very effective at producing quality solutions and can
handle larger instances than previously suggested approaches, which
have been primarily based on two-index formulations. The approach, however, does
not seem as effective at producing quality lower bounds. This is
likely due to the absence of dynamic cut generation.
The next step in this research will be to incorporate the generation
of known classes of valid inequalities into our algorithm. This should
produce an algorithm exhibiting the advantages of both the
branch-and-price and branch-and-cut approaches.

%\pagebreak

\bibliography{lrp.bib}
\bibliographystyle{./styles/spbasic}
%\bibliographystyle{plainnat}
\end{document}
